\chapter{Modular Pipeline Design}
\label{ch:modular-pipeline}

Chapter~\ref{ch:complete-pipeline} gives the logical pipeline order.
Chapter~\ref{ch:architecture} describes the ensemble model architecture.
This chapter describes the software design: config-driven, registry-based, every experiment defined by one YAML file.


%% ──────────────────────────────────────────────
\section{Pipeline Architecture}
\label{sec:pipeline-architecture}
%% ──────────────────────────────────────────────

Figure~\ref{fig:pipeline-plugpoints} shows the four pipeline stages and their plug points.
\textsc{Train} is the only modular stage; features, models, and CV strategies resolve from registries at runtime.

\begin{figure}[htbp]
\centering
\resizebox{\textwidth}{!}{%
\begin{tikzpicture}[
  node distance=0.5cm and 0.8cm,
  >=Stealth,
  stage/.style={draw, rounded corners, minimum height=1.4cm, minimum width=2.6cm,
    align=center, font=\small},
  config/.style={draw, rounded corners=8pt, minimum height=1.2cm, minimum width=2.0cm,
    align=center, font=\small, fill=blue!8, line width=0.8pt},
  trainstage/.style={draw, rounded corners, minimum height=1.6cm, minimum width=3.2cm,
    align=center, font=\small, fill=orange!10, line width=1.0pt},
  registry/.style={draw, rounded corners, minimum height=1.6cm, minimum width=3.2cm,
    align=center, font=\scriptsize, fill=blue!4},
  arrow/.style={-{Stealth[length=2.5mm]}, thick},
  plugarrow/.style={-{Stealth[length=2mm]}, thick, defblue},
  reglabel/.style={font=\scriptsize\bfseries, text=defblue},
]

% ── Pipeline stages (top row) ──
\node[config] (cfg) {\textsc{Config}\\[1pt]{\scriptsize(YAML)}};
\node[stage, fill=green!8, right=1.0cm of cfg] (ingest)
  {\textsc{Ingest}\\[1pt]{\scriptsize ticks $\to$ bars}\\[-1pt]{\scriptsize $\to$ 18 measures}};
\node[trainstage, right=0.8cm of ingest] (train)
  {\textsc{Train}\\[1pt]{\scriptsize features + model}\\[-1pt]{\scriptsize + CV splits}};
\node[stage, fill=green!8, right=0.8cm of train] (eval)
  {\textsc{Evaluate}\\[1pt]{\scriptsize $\QLIKE$, MSE, $R^2$}};
\node[stage, fill=green!5, right=0.8cm of eval, dashed] (compare)
  {\textsc{Compare}\\[1pt]{\scriptsize DM, MCS}\\[-1pt]{\scriptsize\itshape (future)}};

% ── Main flow arrows ──
\draw[arrow] (cfg) -- (ingest);
\draw[arrow] (ingest) -- (train);
\draw[arrow] (train) -- (eval);
\draw[arrow, dashed] (eval) -- (compare);

% ── Registry boxes (below TRAIN) ──
\node[registry, below=1.4cm of train, xshift=-3.8cm] (freg) {%
  \begin{tabular}{@{}l@{}}
  {\bfseries\color{defblue}\textsc{Feature\_Registry}}\\[3pt]
  \texttt{har\_core}\\
  \texttt{asymmetry}\\
  \texttt{options}\\
  \texttt{cross\_asset}\\
  \texttt{...}
  \end{tabular}};

\node[registry, below=1.4cm of train] (mreg) {%
  \begin{tabular}{@{}l@{}}
  {\bfseries\color{defblue}\textsc{Model\_Registry}}\\[3pt]
  \texttt{har}\\
  \texttt{harq}\\
  \texttt{lightgbm}\\
  \texttt{shar}\\
  \texttt{...}
  \end{tabular}};

\node[registry, below=1.4cm of train, xshift=3.8cm] (creg) {%
  \begin{tabular}{@{}l@{}}
  {\bfseries\color{defblue}\textsc{CV Strategies}}\\[3pt]
  \texttt{expanding\_window}\\
  \texttt{rolling\_window}\\
  \texttt{purged\_kfold}\\
  \texttt{blocked\_kfold}
  \end{tabular}};

% ── Plug arrows (registries → TRAIN) ──
\draw[plugarrow] (freg.north) -- ([xshift=-1.4cm]train.south);
\draw[plugarrow] (mreg.north) -- (train.south);
\draw[plugarrow] (creg.north) -- ([xshift=1.4cm]train.south);

\end{tikzpicture}%
}
\caption{Pipeline architecture with plug points.
\textsc{Config} parameterises all stages.
\textsc{Train} resolves features, models, and CV strategies from registries at runtime.
\textsc{Compare} (dashed) is not yet implemented.}
\label{fig:pipeline-plugpoints}
\end{figure}


%% ──────────────────────────────────────────────
\section{Plug Points}
\label{sec:plug-points}
%% ──────────────────────────────────────────────

\subsection{Feature Layers}

The \texttt{feature\_layers} list in the config names the layers to compose.
The pipeline iterates the list, resolves each name from \textsc{Feature\_Registry}, and calls \texttt{.compute()}.
See Table~\ref{tab:feature-matrix} for the full feature inventory.

\subsection{Models}

The \texttt{model.name} field selects the model class from \textsc{Model\_Registry}.
All models satisfy the same \texttt{VolModel} protocol: \texttt{.fit(X, y)} and \texttt{.predict(X)}.
Model-specific hyperparameters go in \texttt{model.params}.
See Chapter~\ref{ch:lightgbm} for LightGBM configuration.

\subsection{Cross-Validation}

The \texttt{cv.method} field selects the splitter: expanding window, rolling window, purged k-fold, or blocked k-fold.
Purge gap is set via \texttt{cv.purge\_gap}.
See Chapter~\ref{ch:evaluation} for evaluation methodology.


%% ──────────────────────────────────────────────
\section{Configuration}
\label{sec:configuration}
%% ──────────────────────────────────────────────

Figure~\ref{fig:config-comparison} shows two experiment configs side by side: a HAR baseline and a full LightGBM run.

\begin{figure}[htbp]
\centering
\begin{minipage}[t]{0.47\textwidth}
\centering
{\small\bfseries HAR baseline}
\begin{verbatim}
  name: baseline_har
  universe: [SPY]
  horizons: [1, 5, 22]
  feature_layers: [har_core]
  model:
    name: har
  cv:
    method: expanding_window
    purge_gap: 5
\end{verbatim}
\end{minipage}%
\hfill
\begin{minipage}[t]{0.47\textwidth}
\centering
{\small\bfseries LightGBM with options features}
\begin{verbatim}
  name: lgbm_full
  universe: [SPY, AAPL, MSFT, NVDA]
  horizons: [1, 5, 22]
  feature_layers:
    - har_core
    - asymmetry
    - options
    - cross_asset
  model:
    name: lightgbm
    params: {num_leaves: 31}
  cv:
    method: purged_kfold
    purge_gap: 22
\end{verbatim}
\end{minipage}
\caption{Two experiments, same pipeline.
Left: HAR baseline with one feature layer.
Right: LightGBM with four feature layers and stricter cross-validation.}
\label{fig:config-comparison}
\end{figure}


%% ──────────────────────────────────────────────
%% Boxes
%% ──────────────────────────────────────────────

\begin{keyidea}[One YAML, One Experiment]
The entire experiment (universe, horizons, features, model, CV strategy) is defined by a single config file.
No code changes needed to run a different experiment.
\end{keyidea}

\begin{warning}[Registries Require Import]
Feature layers and models register via decorators at import time.
A new layer or model that is never imported will not appear in the registry.
The package \texttt{\_\_init\_\_.py} triggers all registration imports.
\end{warning}
